\section{AHSME 1956}
        \begin{enumerate}
        	\item (\textit{AHSME 1956}) The value of $x + x(x^x)$ when $x = 2$ is: 


 $\textbf{(A)}\ 10 \qquad\textbf{(B)}\ 16 \qquad\textbf{(C)}\ 18 \qquad\textbf{(D)}\ 36 \qquad\textbf{(E)}\ 64$


	\item (\textit{AHSME 1956}) Mr. Jones sold two pipes at $\textdollar{ 1.20}$ each. Based on the cost, his profit on one was


 $20$\% and his loss on the other was $20$\%. 
On the sale of the pipes, he: 


 $\textbf{(A)}\ \text{broke even}\qquad \textbf{(B)}\ \text{lost }4\text{ cents} \qquad\textbf{(C)}\ \text{gained }4\text{ cents}\qquad \\ \textbf{(D)}\ \text{lost }10\text{ cents}\qquad \textbf{(E)}\ \text{gained }10\text{ cents}$


	\item (\textit{AHSME 1956}) The distance light travels in one year is approximately $5,870,000,000,000$ miles. The distance light travels in $100$ years is: 


 $\textbf{(A)}\ 587 * 10^8\text{ miles}\qquad \textbf{(B)}\ 587 * 10^{10}\text{ miles}\qquad \textbf{(C)}\ 587*10^{-10}\text{ miles} \\ \textbf{(D)}\ 587 * 10^{12} \text{ miles} \qquad \textbf{(E)}\ 587* 10^{ - 12} \text{ miles}$


	\item (\textit{AHSME 1956}) A man has $\textdollar{10,000 }$ to invest. He invests $\textdollar{4000}$ at 5\% and $\textdollar{3500}$ at 4\%. 
In order to have a yearly income of $\textdollar{500}$, he must invest the remainder at: 


 $\textbf{(A)}\ 6\%\qquad\textbf{(B)}\ 6.1\%\qquad\textbf{(C)}\ 6.2\%\qquad\textbf{(D)}\ 6.3\%\qquad\textbf{(E)}\ 6.4\%$


	\item (\textit{AHSME 1956}) A nickel is placed on a table. The number of nickels which can be placed around it, each tangent to it and to two others is: 


 $\textbf{(A)}\ 4 \qquad\textbf{(B)}\ 5 \qquad\textbf{(C)}\ 6 \qquad\textbf{(D)}\ 8 \qquad\textbf{(E)}\ 12$


	\item (\textit{AHSME 1956}) In a group of cows and chickens, the number of legs was 14 more than twice the number of heads. The number of cows was: 


 $\textbf{(A)}\ 5 \qquad\textbf{(B)}\ 7 \qquad\textbf{(C)}\ 10 \qquad\textbf{(D)}\ 12 \qquad\textbf{(E)}\ 14$


	\item (\textit{AHSME 1956}) The roots of the equation $ax^2 + bx + c = 0$ will be reciprocal if: 


 $\textbf{(A)}\ a = b \qquad\textbf{(B)}\ a = bc \qquad\textbf{(C)}\ c = a \qquad\textbf{(D)}\ c = b \qquad\textbf{(E)}\ c = ab$


	\item (\textit{AHSME 1956}) If $8\cdot2^x = 5^{y + 8}$, then when $y = - 8,x =$
$\textbf{(A)}\ - 4 \qquad\textbf{(B)}\ - 3 \qquad\textbf{(C)}\ 0 \qquad\textbf{(D)}\ 4 \qquad\textbf{(E)}\ 8$


	\item (\textit{AHSME 1956}) When you simplify $\left[ \sqrt [3]{\sqrt [6]{a^9}} \right]^4\left[ \sqrt [6]{\sqrt [3]{a^9}} \right]^4$, the result is: 


 $\textbf{(A)}\ a^{16} \qquad\textbf{(B)}\ a^{12} \qquad\textbf{(C)}\ a^8 \qquad\textbf{(D)}\ a^4 \qquad\textbf{(E)}\ a^2$


	\item (\textit{AHSME 1956}) A circle of radius $10$ inches has its center at the vertex $C$ of an equilateral $\triangle ABC$ and passes through the other two vertices. 
The side $AC$ extended through $C$ intersects the circle at $D$. The number of degrees of $\angle ADB$ is: 


 $\textbf{(A)}\ 15^{\circ} \qquad \textbf{(B)}\ 30^{\circ} \qquad \textbf{(C)}\ 60^{\circ} \qquad \textbf{(D)}\ 90^{\circ} \qquad \textbf{(E)}\ 120^{\circ}$


	\item (\textit{AHSME 1956}) The expression $1 - \frac {1}{1 + \sqrt {3}} + \frac {1}{1 - \sqrt {3}}$ equals: 


 $\textbf{(A)}\ 1-\sqrt{3}\qquad \textbf{(B)}\ 1\qquad \textbf{(C)}\ -\sqrt{3}\qquad \textbf{(D)}\ \sqrt{3}\qquad \textbf{(E)}\ 1+\sqrt{3}$


	\item (\textit{AHSME 1956}) If $x^{ - 1} - 1$ is divided by $x - 1$ the quotient is: 


 $\textbf{(A)}\ 1\qquad \textbf{(B)}\ \frac{1}{x-1}\qquad \textbf{(C)}\ \frac{-1}{x-1}\qquad \textbf{(D)}\ \frac{1}{x}\qquad \textbf{(E)}\ -\frac{1}{x}$


	\item (\textit{AHSME 1956}) Given two positive integers $x$ and $y$ with $x < y$. The percent that $x$ is less than $y$ is: 


 $\textbf{(A)}\ \frac{100(y-x)}{x}\qquad \textbf{(B)}\ \frac{100(x-y)}{x}\qquad \textbf{(C)}\ \frac{100(y-x)}{y}\qquad \\ \textbf{(D)}\ 100(y-x) \textbf{(E)}\ 100(x - y)$


	\item (\textit{AHSME 1956}) The points $A,B,C$ are on a circle $O$. The tangent line at $A$ and the secant $BC$ intersect at $P, B$ lying between $C$ and $P$. 
If $\overline{BC} = 20$ and $\overline{PA} = 10\sqrt {3}$, then $\overline{PB}$ equals: 


 $\textbf{(A)}\ 5 \qquad\textbf{(B)}\ 10 \qquad\textbf{(C)}\ 10\sqrt {3} \qquad\textbf{(D)}\ 20 \qquad\textbf{(E)}\ 30$


	\item (\textit{AHSME 1956}) The root(s) of $\frac {15}{x^2 - 4} - \frac {2}{x - 2} = 1$ is (are): 


 $\textbf{(A)}\ -5\text{ and }3\qquad \textbf{(B)}\ \pm 2\qquad \textbf{(C)}\ 2\text{ only}\qquad \textbf{(D)}\ -3\text{ and }5\qquad \textbf{(E)}\ 3\text{ only}$


	\item (\textit{AHSME 1956}) The sum of three numbers is $98$. The ratio of the first to the second is $\frac {2}{3}$, 
and the ratio of the second to the third is $\frac {5}{8}$. The second number is: 


 $\textbf{(A)}\ 15 \qquad\textbf{(B)}\ 20 \qquad\textbf{(C)}\ 30 \qquad\textbf{(D)}\ 32 \qquad\textbf{(E)}\ 33$


	\item (\textit{AHSME 1956}) The fraction $\frac {5x - 11}{2x^2 + x - 6}$ was obtained by adding the two fractions $\frac {A}{x + 2}$ and 
$\frac {B}{2x - 3}$. The values of $A$ and $B$ must be, respectively: 


 $\textbf{(A)}\ 5x,-11\qquad\textbf{(B)}\ -11,5x\qquad\textbf{(C)}\ -1,3\qquad\textbf{(D)}\ 3,-1\qquad\textbf{(E)}\ 5,-11$


	\item (\textit{AHSME 1956}) If $10^{2y} = 25$, then $10^{ - y}$ equals: 


 $\textbf{(A)}\ -\frac{1}{5}\qquad \textbf{(B)}\ \frac{1}{625}\qquad \textbf{(C)}\ \frac{1}{50}\qquad \textbf{(D)}\ \frac{1}{25}\qquad \textbf{(E)}\ \frac{1}{5}$


	\item (\textit{AHSME 1956}) Two candles of the same height are lighted at the same time. The first is consumed in $4$ hours and the second in $3$ hours. 
Assuming that each candle burns at a constant rate, in how many hours after being lighted was the first candle twice the height of the second? 


 $\textbf{(A)}\ \frac{3}{4}\qquad\textbf{(B)}\ 1\frac{1}{2}\qquad\textbf{(C)}\ 2\qquad\textbf{(D)}\ 2\frac{2}{5}\qquad\textbf{(E)}\ 2\frac{1}{2}$


	\item (\textit{AHSME 1956}) If $(0.2)^x = 2$ and $\log 2 = 0.3010$, then the value of $x$ to the nearest tenth is: 


 $\textbf{(A)}\ - 10.0 \qquad\textbf{(B)}\ - 0.5 \qquad\textbf{(C)}\ - 0.4 \qquad\textbf{(D)}\ - 0.2 \qquad\textbf{(E)}\ 10.0$


	\item (\textit{AHSME 1956}) If each of two intersecting lines intersects a hyperbola and neither line is tangent to the hyperbola, 
then the possible number of points of intersection with the hyperbola is: 


 $\textbf{(A)}\ 2\qquad \textbf{(B)}\ 2\text{ or }3\qquad \textbf{(C)}\ 2\text{ or }4\qquad \textbf{(D)}\ 3\text{ or }4\qquad \textbf{(E)}\ 2,3,\text{ or }4$


	\item (\textit{AHSME 1956}) Jones covered a distance of $50$ miles on his first trip. On a later trip he traveled $300$ miles while going three times as fast. 
His new time compared with the old time was: 


 $\textbf{(A)}\ \text{three times as much} \qquad \textbf{(B)}\ \text{twice as much} \qquad \textbf{(C)}\ \text{the same} \\ \textbf{(D)}\ \text{half as much} \qquad \textbf{(E)}\ \text{a third as much}$


	\item (\textit{AHSME 1956}) About the equation $ax^2 - 2x\sqrt {2} + c = 0$, with $a$ and $c$ real constants, 
we are told that the discriminant is zero. The roots are necessarily: 


 $\textbf{(A)}\ \text{equal and integral}\qquad \textbf{(B)}\ \text{equal and rational}\qquad \textbf{(C)}\ \text{equal and real} \\ \textbf{(D)}\ \text{equal and irrational} \qquad \textbf{(E)}\ \text{equal and imaginary}$


	\item (\textit{AHSME 1956}) In the figure $\overline{AB} = \overline{AC}$, $\angle BAD = 30^{\circ}$, and $\overline{AE} = \overline{AD}$. 


 
                    \begin{figure}[H]
                        \centering
                        \includegraphics[width=0.5\textwidth]{Problem/AHSME_I_1956_24_Problem_Figure_1.png}
                    \end{figure}
                    
The angle CDE equals:


 $\textbf{(A)}\ 7\frac{1}{2}^{\circ}\qquad \textbf{(B)}\ 10^{\circ}\qquad \textbf{(C)}\ 12\frac{1}{2}^{\circ}\qquad \textbf{(D)}\ 15^{\circ}\qquad \textbf{(E)}\ 20^{\circ}$


	\item (\textit{AHSME 1956}) The sum of all numbers of the form $2k + 1$, where $k$ takes on integral values from $1$ to $n$ is: 


 $\textbf{(A)}\ n^2\qquad\textbf{(B)}\ n(n+1)\qquad\textbf{(C)}\ n(n+2)\qquad\textbf{(D)}\ (n+1)^2\qquad\textbf{(E)}\ (n+1)(n+2)$


	\item (\textit{AHSME 1956}) Which one of the following combinations of given parts does not determine the indicated triangle? 


 $\textbf{(A)}\ \text{base angle and vertex angle; isosceles triangle} \\ \textbf{(B)}\ \text{vertex angle and the base; isosceles triangle} \\ \textbf{(C)}\ \text{the radius of the circumscribed circle; equilateral triangle} \\ \textbf{(D)}\ \text{one arm and the radius of the inscribed circle; right triangle} \\ \textbf{(E)}\ \text{two angles and a side opposite one of them; scalene triangle}$


	\item (\textit{AHSME 1956}) If an angle of a triangle remains unchanged but each of its two including sides is doubled, then the area is multiplied by: 


 $\textbf{(A)}\ 2 \qquad\textbf{(B)}\ 3 \qquad\textbf{(C)}\ 4 \qquad\textbf{(D)}\ 6 \qquad\textbf{(E)}\ \text{more than }6$


	\item (\textit{AHSME 1956}) Mr. J left his entire estate to his wife, his daughter, his son, and the cook. 
His daughter and son got half the estate, sharing in the ratio of $4$ to $3$. 
His wife got twice as much as the son. If the cook received a bequest of $\textdollar{500}$, then the entire estate was: 


 $\textbf{(A)}\ \textdollar{3500}\qquad \textbf{(B)}\ \textdollar{5500}\qquad \textbf{(C)}\ \textdollar{6500}\qquad \textbf{(D)}\ \textdollar{7000}\qquad \textbf{(E)}\ \textdollar{7500}$


	\item (\textit{AHSME 1956}) The points of intersection of $xy = 12$ and $x^2 + y^2 = 25$ are joined in succession. The resulting figure is: 


 $\textbf{(A)}\ \text{a straight line}\qquad \textbf{(B)}\ \text{an equilateral triangle}\qquad \textbf{(C)}\ \text{a parallelogram} \\ \textbf{(D)}\ \text{a rectangle} \qquad \textbf{(E)}\ \text{a square}$


	\item (\textit{AHSME 1956}) If the altitude of an equilateral triangle is $\sqrt {6}$, then the area is: 


 $\textbf{(A)}\ 2\sqrt{2}\qquad\textbf{(B)}\ 2\sqrt{3}\qquad\textbf{(C)}\ 3\sqrt{3}\qquad\textbf{(D)}\ 6\sqrt{2}\qquad\textbf{(E)}\ 12$


	\item (\textit{AHSME 1956}) In our number system the base is ten. If the base were changed to four you would count as follows: 
$1,2,3,10,11,12,13,20,21,22,23,30,\ldots$ The twentieth number would be: 


 $\textbf{(A)}\ 20 \qquad\textbf{(B)}\ 38 \qquad\textbf{(C)}\ 44 \qquad\textbf{(D)}\ 104 \qquad\textbf{(E)}\ 110$


	\item (\textit{AHSME 1956}) George and Henry started a race from opposite ends of the pool. After a minute and a half, 
they passed each other in the center of the pool. If they lost no time in turning and 
maintained their respective speeds, how many minutes after starting did they pass each other the second time? 


 $\textbf{(A)}\ 3\qquad \textbf{(B)}\ 4\frac{1}{2}\qquad \textbf{(C)}\ 6\qquad \textbf{(D)}\ 7\frac{1}{2}\qquad \textbf{(E)}\ 9$


	\item (\textit{AHSME 1956}) The number $\sqrt {2}$ is equal to: 


 $\textbf{(A)}\ \text{a rational fraction} \qquad \textbf{(B)}\ \text{a finite decimal} \qquad \textbf{(C)}\ 1.41421 \\ \textbf{(D)}\ \text{an infinite repeating decimal} \qquad \textbf{(E)}\ \text{an infinite non - repeating decimal}$


	\item (\textit{AHSME 1956}) If $n$ is any whole number, $n^2(n^2 - 1)$ is always divisible by 


 $\textbf{(A)}\ 12\qquad \textbf{(B)}\ 24\qquad \textbf{(C)}\ \text{any multiple of }12\qquad \textbf{(D)}\ 12-n\qquad \textbf{(E)}\ 12\text{ and }24$


	\item (\textit{AHSME 1956}) A rhombus is formed by two radii and two chords of a circle whose radius is $16$ feet. The area of the rhombus in square feet is: 


 $\textbf{(A)}\ 128\qquad \textbf{(B)}\ 128\sqrt{3}\qquad \textbf{(C)}\ 256\qquad \textbf{(D)}\ 512\qquad \textbf{(E)}\ 512\sqrt{3}$


	\item (\textit{AHSME 1956}) If the sum $1 + 2 + 3 + \cdots + K$ is a perfect square $N^2$ and if $N$ is less than $100$, then the possible values for $K$ are: 


 $\textbf{(A)}\ \text{only }1\qquad \textbf{(B)}\ 1\text{ and }8\qquad \textbf{(C)}\ \text{only }8\qquad \textbf{(D)}\ 8\text{ and }49\qquad \textbf{(E)}\ 1,8,\text{ and }49$


	\item (\textit{AHSME 1956}) On a map whose scale is $400$ miles to an inch and a half, a certain estate is represented by a 
rhombus having a $60^{\circ}$ angle. The diagonal opposite $60^{\circ}$ is $\frac {3}{16}$ in. The area of the estate in square miles is: 


 $\textbf{(A)}\ \frac{2500}{\sqrt{3}}\qquad \textbf{(B)}\ \frac{1250}{\sqrt{3}}\qquad \textbf{(C)}\ 1250\qquad \textbf{(D)}\ \frac{5625\sqrt{3}}{2}\qquad \textbf{(E)}\ 1250\sqrt{3}$


	\item (\textit{AHSME 1956}) In a right triangle with sides $a$ and $b$, and hypotenuse $c$, the altitude drawn on the hypotenuse is $x$. Then: 


 $\textbf{(A)}\ ab = x^2 \qquad \textbf{(B)}\ \frac {1}{a} + \frac {1}{b} = \frac {1}{x} \qquad \textbf{(C)}\ a^2 + b^2 = 2x^2 \\ \textbf{(D)}\ \frac {1}{x^2} = \frac {1}{a^2} + \frac {1}{b^2} \qquad \textbf{(E)}\ \frac {1}{x} = \frac {b}{a}$


	\item (\textit{AHSME 1956}) The hypotenuse $c$ and one arm $a$ of a right triangle are consecutive integers. The square of the second arm is: 


 $\textbf{(A)}\ ca\qquad \textbf{(B)}\ \frac{c}{a}\qquad \textbf{(C)}\ c+a\qquad \textbf{(D)}\ c-a\qquad  \textbf{(E)}\ \text{none of these}$


	\item (\textit{AHSME 1956}) If $V = gt + V_0$ and $S = \frac {1}{2}gt^2 + V_0t$, then $t$ equals: 


 $\textbf{(A)}\ \frac{2S}{V+V_0}\qquad \textbf{(B)}\ \frac{2S}{V-V_0}\qquad \textbf{(C)}\ \frac{2S}{V_0-V}\qquad \textbf{(D)}\ \frac{2S}{V}\qquad \textbf{(E)}\ 2S-V$


	\item (\textit{AHSME 1956}) The equation $3y^2 + y + 4 = 2(6x^2 + y + 2)$ where $y = 2x$ is satisfied by: 


 $\textbf{(A)}\ \text{no value of }x \qquad \textbf{(B)}\ \text{all values of }x \qquad \textbf{(C)}\ x = 0\text{ only} \\ \textbf{(D)}\ \text{all integral values of }x\text{ only} \qquad \textbf{(E)}\ \text{all rational values of }x\text{ only}$


	\item (\textit{AHSME 1956}) The equation $\sqrt {x + 4} - \sqrt {x - 3} + 1 = 0$ has: 


 $\textbf{(A)}\ \text{no root} \qquad  \textbf{(B)}\ \text{one real root} \\ \textbf{(C)}\ \text{one real root and one imaginary root} \\ \textbf{(D)}\ \text{two imaginary roots} \qquad \qquad\textbf{(E)}\ \text{two real roots}$


	\item (\textit{AHSME 1956}) The number of scalene triangles having all sides of integral lengths, and perimeter less than $13$ is: 


 $\textbf{(A)}\ 1 \qquad\textbf{(B)}\ 2 \qquad\textbf{(C)}\ 3 \qquad\textbf{(D)}\ 4 \qquad\textbf{(E)}\ 18$


	\item (\textit{AHSME 1956}) If $x < a < 0$ means that $x$ and $a$ are numbers such that $x$ is less than $a$ and $a$ is less than zero, then: 


 $\textbf{(A)}\ x^2 < ax < 0 \qquad \textbf{(B)}\ x^2 > ax > a^2 \qquad \textbf{(C)}\ x^2 < a^2 < 0 \\ \textbf{(D)}\ x^2 > ax\text{ but }ax < 0 \qquad \textbf{(E)}\ x^2 > a^2\text{ but }a^2 < 0$


	\item (\textit{AHSME 1956}) A wheel with a rubber tire has an outside diameter of $25$ in. 
When the radius has been decreased a quarter of an inch, the number of revolutions in one mile will: 


 $\text{(A)}\ \text{be increased about }2\% \qquad \\ \text{(B)}\ \text{be increased about }1\%  \\ \text{(C)}\ \text{be increased about }20\%\qquad \\ \text{(D)}\ \text{be increased about }\frac{1}{2}\%\qquad \\ \text{(E)}\ \text{remain the same}$


	\item (\textit{AHSME 1956}) For the equation $\frac {1 + x}{1 - x} = \frac {N + 1}{N}$ to be true where $N$ is positive, $x$ can have: 


 $\textbf{(A)}\ \text{any positive value less than }1 \qquad \\ \textbf{(B)}\ \text{any value less than }1 \\ \textbf{(C)}\ \text{the value zero only}\qquad \\ \textbf{(D)}\ \text{any non-negative value}\qquad \\ \textbf{(E)}\ \text{any value}$


	\item (\textit{AHSME 1956}) An engineer said he could finish a highway section in $3$ days with his present supply of a certain type of machine. 
However, with $3$ more of these machines the job could be done in $2$ days. If the machines all work at the same rate, 
how many days would it take to do the job with one machine? 


 $\textbf{(A)}\ 6 \qquad \textbf{(B)}\ 12 \qquad \textbf{(C)}\ 15 \qquad \textbf{(D)}\ 18 \qquad \textbf{(E)}\ 36$


	\item (\textit{AHSME 1956}) If $p$ is a positive integer, then $\frac {3p + 25}{2p - 5}$ can be a positive integer, if and only if $p$ is: 


 $\textbf{(A)}\ \text{at least }3\qquad \textbf{(B)}\ \text{at least }3\text{ and no more than }35\qquad \\ \textbf{(C)}\ \text{no more than }35  \textbf{(D)}\ \text{equal to }35 \qquad \textbf{(E)}\ \text{equal to }3\text{ or }35$


	\item (\textit{AHSME 1956}) Triangle $PAB$ is formed by three tangents to circle $O$ and $\angle APB = 40^{\circ}$; then $\angle AOB$ equals: 


 $\textbf{(A)}\ 45^{\circ}\qquad \textbf{(B)}\ 50^{\circ}\qquad \textbf{(C)}\ 55^{\circ}\qquad \textbf{(D)}\ 60^{\circ}\qquad \textbf{(E)}\ 70^{\circ}$


	\item (\textit{AHSME 1956}) In $\triangle ABC, \overline{CA} = \overline{CB}$. On $CB$ square $BCDE$ is constructed away from the triangle. 
If $x$ is the number of degrees in $\angle DAB$, then 


 $\textbf{(A)}\ x\text{ depends upon }\triangle ABC \qquad \textbf{(B)}\ x\text{ is independent of the triangle} \\ \textbf{(C)}\ x\text{ may equal }\angle CAD \qquad \\ \textbf{(D)}\ x\text{ can never equal }\angle CAB \\ \textbf{(E)}\ x\text{ is greater than }45^{\circ}\text{ but less than }90^{\circ}$


\end{enumerate}